%%%%%%%%%%%%%%%%%%%%%%%%%%%%%%%%%%%%%%%%%%%%%%%%%%%%%%%%%%%%%%%%%%%%%%%%
\chapter{Applicazioni della CP alla complessità aritmetica}
%%%%%%%%%%%%%%%%%%%%%%%%%%%%%%%%%%%%%%%%%%%%%%%%%%%%%%%%%%%%%%%%%%%%%%%%

Introduciamo una quantità per misurare la complessità di una moltiplicazione tra matrici:
\begin{definition}
  L'\emph{esponente} $ \omega $ di una moltiplicazione matriciale è il numero
\begin{equation*}
  \omega = \inf\{h \in \R \mid \text{ la moltiplicazione tra due matrici }  n\times n  \text{ ha costo } O(n^h) \text{ operazioni aritmetiche.} \}.
\end{equation*}
\end{definition}
L'algoritmo di Strassen mostra che $ \omega \le \log_2(7) < 2.81 $, e si ha $ \omega \ge 2 $ dal
costo $ O(n^2) $ per il calcolo di addizioni tra matrici. Trovare $ \omega $ è un problema centrale
nella teoria della complessità, e si congettura che $ \omega = 2 $, cioè che il costo della moltiplicazioni possa essere
asintoticamente pari a quello
delle moltiplicazioni. Il record attuale è $ \omega \le 2.371339 $~\cite{alman2024}. Dopo Strassen,
nel 1978 i miglioramenti significativi sono stati $ \omega \le 2.7962 $ ad opera di
\cite{pan78}, poi $ \omega \le 2.7799 $ \cite{bini1980relations} in 1979, poi $ \omega \le 2.55 $ \cite{sch81}, poi $
\omega \le  2.48 $ nel 1887 \cite{str87}, e poi $ \omega \le 2.38 $ \cite{cw90} nel 1989, che potrebbe aver fatto
pensare che nel 1990 si sarebbe trovato il bound. Dopo tale miglioramento, l'approssimazione non è
più scesa sotto $ \omega < 2.373 $ \cite{wil11,gall14,stothers2010complexity}.

\begin{remark}
  Trovare il minimo esponente è un problema di interesse principalmente della teoria della
  complessità. Nella maggior parte della applicazioni, l'algoritmo di Strassen è quello ottimale, si
  veda la discussione in \cite{dalberto2023strassensmatrixmultiplicationalgorithm}.
\end{remark}

\section{Rappresentare forme bilineari tramite tensori}
\label{sec:tensor_bil}

\todo{Rename bilinear form with f}

Conoscere il rango CP di tensori che rappresentano applicazioni multilineari può rivelare
informazioni sulla complessità asintotica e fornire algoritmi con complessità ottimale per la loro valutazione. 
Proviamo a formulare il problema. Dall'algebra lineare sappiamo che data una forma bilineare $ b : \R^m \times \R^n \to \R $, esiste una matrice $ M \in \R^{m \times n} $ tale che $
b(x,y) = x^{\top}My$ per ogni $ (x,y) \in \R^m \times \R^n $.
Notiamo che un insieme di $ p $ forme bilineari, o equivalentemente una funzione \todo{È chiaro se
si chiama ancora b?}bilineare $
b = \begin{bmatrix} b_1 \\ \vdots \\ b_p \end{bmatrix} :  \R^m \times \R^n  \to \R^p$ si può
rappresentare tramite un tensore $ \Tn{t} \in \R^{m \times n\times p} $ tale che le sue slice
frontali rappresentino ognuna delle $ p $ forme bilineari, ossia che
\begin{equation}
  \label{eq:b_k_bil}
  b_k(x,y) = x^{\top}\Tn{t}(:,:,k)y = \sum_{i=1}^m\sum_{j=1}^n \Tn{t}_{i,j,k}x_iy_j, \text{ per ogni
  } k = 1, \ldots, p
.\end{equation}
Da cui, 
\begin{equation}
  \label{eq:b_tuv}
 b(x,y) = \begin{bmatrix} b_1(x,y) \\ \vdots \\ b_p(x,y) \end{bmatrix} = \Tn{t} \ttm[1] x^{\top}
 \ttm[2] y^{\top}.
\end{equation}
\begin{remark}
  Questa procedura si estende anche al caso in cui le entrate di $ x $ e $ y $ siano oggetti
  algebrici più complessi di scalari, ad esempio vettori o matrici.
\end{remark}
\begin{remark}[Active multiplications]
  Scrivere una forma bilineare come nell'Equazione~\ref{eq:b_k_bil} evidenzia il numero di
  moltiplicazioni necessarie fra gli elementi di $ x $ e $ y $. Queste sono dette \emph{active
  multiplications} e corrispondono al numero di elementi non nulli di $ \Tn{t} $. Poiché ci aspettiamo che
  il costo delle active multiplications sia quello che determini la complessità di valutare $ b(x,y)
  $, cerchiamo di minimizzarne il numero cercando fattorizzazioni CP di $ \Tn{t} $.
\end{remark}
Data una decomposizione $ \Tn{t} = \dsqr{U,V,W} $ di rango $ r $, dalla Definizione~\ref{def:CP} abbiamo che $ \Tn{t} = \sum_{l=1}^r u_l \topp v_l \topp w_l $, o in componenti, $ \Tn{t}(i,j,k) = \sum_{l=1}^r u_{il}v_{jl}w_{kl} $. Inserendo questa relazione nella Equazione~\ref{eq:b_k_bil} troviamo,

\begin{equation}
  \label{eq:b_active}
  \begin{aligned}
    b_k(x,y) &= \left( \left( \sum_{l=1}^r u_l \topp v_l \topp w_l \right) \ttm[1] x^{\top} \ttm[2]
  y^{\top}  \right)_k  
             = \left( \sum_{l=1}^r \left( u_l \topp v_l \topp w_l  \right) \ttm[1] x^{\top} \ttm[2]
    y^{\top} \right)_k  \\ 
             &=
  \sum_{i=1}^m \sum_{j=1}^n \sum_{l=1}^r u_{il}v_{jl}w_{kl} x_i y_j = \sum_{l=1}^r \left(\sum_{i=1}^m 
   u_{il} x_i \right) \cdot \left( \sum_{j=1}^n v_{jl} y_j  \right)  w_{kl} = \sum_{l=1}^r
   (u_l^{\top}x) \cdot (v_l^{\top}y)w_{kl},
  \end{aligned}
\end{equation}
per ogni $ k = 1, \ldots, p $.
Mettendo insieme le Equazioni \ref{eq:b_tuv} e \ref{eq:b_active} troviamo
\begin{equation*}
  b(x,y) = \sum_{l=1}^r (u_l^{\top}x)\cdot(v_l^{\top}y)w_l.
\end{equation*}

Quindi valutare $ b $ richiede esattamente $ r $ active multiplications, che evidenziamo con $
(\cdot) $.
\todo{Espandere discussione dall'articolo}



\section{L'algoritmo di Strassen}

Possiamo rappresentare la moltiplicazione tra matrici come una funzione bilineare, quindi come un
tensore. 

\todo{Brief motivation}

\begin{notation}
  Indichiamo con $ \angb{m,n,p} $ il prodotto tra due matrici di taglia $ m \times n $ e $ n \times
  p $.
\end{notation}

\subsection*{Il caso naive}

Trattiamo il caso di matrici quadrate con taglie potenze di due (possiamo farlo a meno di passare a
una sottomatrice di taglia pari). In questo caso possiamo
partizionare le matrici in delle sottomatrici a blocchi $ 2\times 2 $:
\begin{equation*}
  \begin{bmatrix}
    C_{11} & C_{12} \\
    C_{21} & C_{22}
  \end{bmatrix} =
  \begin{bmatrix}
    A_{11} & A_{12} \\
    A_{21} & A_{22}
  \end{bmatrix} \cdot 
  \begin{bmatrix}
    B_{11} & B_{12} \\
    B_{21} & B_{22}
  \end{bmatrix} =
  \begin{bmatrix}
    A_{11}B_{11} + A_{12}B_{21} & A_{11}B_{12} + A_{12}B_{22} \\
    A_{21}B_{11} + A_{22}B_{21} & A_{21}B_{12} + A_{22}B_{22}
  \end{bmatrix}.
\end{equation*}
L'algoritmo naive procede combinando un insieme di otto moltiplicazioni matriciali con quattro
addizioni:
% Intermediate Matrices
% Intermediate Matrices
\[
\begin{aligned}
M_1 &= A_{11} \cdot B_{11} &\quad M_2 &= A_{12} \cdot B_{21} &\quad M_3 &= A_{11} \cdot B_{12} \\
M_4 &= A_{12} \cdot B_{22} &\quad M_5 &= A_{21} \cdot B_{11} &\quad M_6 &= A_{22} \cdot B_{21} \\
M_7 &= A_{21} \cdot B_{12} &\quad M_8 &= A_{22} \cdot B_{22} & &
\end{aligned}
\]

% Resulting Matrices
\[
\begin{aligned}
C_{11} &= M_1 + M_2 &\qquad C_{12} &= M_3 + M_4 \\
C_{21} &= M_5 + M_6 &\qquad C_{22} &= M_7 + M_8
\end{aligned}
\]

\section{Il problema del border rank}
\label{sec:border_rank}

\todo{Introduci e motiva il problema partendo dagli algoritmi APA (vedi Landsbeg)}

\begin{example}
  Consideriamo dei vettori $ a_1,a_2 \in \R^m$, $b_1,b_2 \in \R^n$ e $c_1,c_2\in \R^p $ vettori
  linearmente indipendenti.  Definiamo il tensore 
  \begin{equation}
    \Tn{X} = a_1 \topp  b_1 \topp c_2 + a_1 \topp b_2 \topp c_1 + a_2 \topp b_1 \topp c_1,
  \end{equation}
  e notiamo che $ \rk(\Tn{x}) = 3 $ dato che non può
  essere scomposto come somma di tensori più semplici.
  Consideriamo la successioni di tensori 
  \begin{equation}
  \label{eq:rank_counterex}
  \Tn{y}_\varepsilon = \frac{1}{\varepsilon}\left( a_1 +
  \varepsilon a_2 \right) \topp (b_1 + \varepsilon b_2) \topp (c_1 + \varepsilon c_2) -
  \frac{1}{\varepsilon} a_1 \topp b_1 \topp c_1, \text{ con } \varepsilon \in \R,
  \end{equation}
  che essendo combinazione lineare di due tensori di rango uno, ha rango 2.
  Vediamo che $ \lim_{\varepsilon \to 0} \Tn{y}_\varepsilon = \Tn{x} $, da cui segue che l'insieme dei tensori di rango 3 non è chiuso.
  Espandendo i termini dell'Equazione \ref{eq:rank_counterex} tramite la proprietà distributiva del prodotto esterno,
  \begin{equation}
    \begin{aligned}
      \Tn{y}_\varepsilon &=  \frac{1}{\varepsilon} a_1 \topp b_1 \topp c_2 + a_1 \topp b_1 \topp c_2 + a_1 \topp b_2 \topp c_1
      +a_2 \topp b_1 \topp c_1  \\ &+ \varepsilon \left(a_2 \topp b_2 \topp c_1 + a_1 \topp b_1 \topp c_2 +
      a_1 \topp b_2 \topp c_2 + a_2 \topp b_2 \topp c_2 \right) - \frac{1}{\varepsilon} a_1 \topp
      b_1 \topp c_1 \\
                         &= \Tn{x} + \varepsilon\left(a_2 \topp b_2 \topp c_1 + a_1 \topp b_1 \topp c_2 +
      a_1 \topp b_2 \topp c_2 + a_2 \topp b_2 \topp c_2 \right)
    \end{aligned}
  \end{equation}
  Quindi $ \Tn{y}_\varepsilon - \Tn{x} \to 0 $.
\end{example}

Questo motiva la seguente definizione:

\begin{definition}[Border rank]
  Dato un tensore $ \Tn{x} \in \Rmsiz $, il suo border rank è la quantità
  \[
  \brk(\Tn{x}) = \min \left\{ r \in \N \mid \exists \Tn{y} \in \Rmsiz! \text{ t.c }
  \rk(\Tn{y})=r, \, \|\Tn{x} - \Tn{y}\|  < \varepsilon \right\}
  .\] 
\end{definition}

